%% ============================================================
%%  CHAPTER 2 --- LITERATURE REVIEW & BACKGROUND
%% ============================================================
\chapter{Literature Review and Background}
\label{ch:literature}

\section{Overview}

This chapter surveys the state of the art across the four principal knowledge domains that underpin \ACRQA{}: static application security testing, AI-assisted code analysis, retrieval-augmented generation, and cryptographic inventory systems. For each domain, foundational concepts are introduced, key prior work is summarised, and the gaps that motivate this project are identified.

% -----------------------------------------------------------------
\section{Static Application Security Testing (SAST)}
\label{sec:sast_background}

\subsection{Definition and Scope}

Static Application Security Testing is the practice of analysing source code, bytecode, or binary representations of software to identify security vulnerabilities without executing the program. Unlike Dynamic Application Security Testing (DAST), which probes a running system, SAST operates on artefacts available during development --- making it naturally suited for integration into CI/CD pipelines \cite{chess2004static}.

SAST tools typically operate using one or more of the following underlying techniques \cite{chess2004static,livshits2005finding}:

\begin{itemize}
  \item \textbf{Pattern matching / taint analysis:} Rules expressed as regular expressions or abstract syntax tree (AST) patterns detect known vulnerability signatures. Tools like Semgrep and PMD operate primarily in this mode.

  \item \textbf{Data flow analysis:} Tracks how untrusted input propagates through program execution paths to sensitive sinks (e.g., database queries, shell commands) \cite{newsome2005dynamic,enck2014taintdroid}. Finds taint-propagation vulnerabilities such as SQL injection and command injection \cite{halfond2006classification}.

  \item \textbf{Control flow analysis:} Models all possible execution paths through a program to identify paths that reach insecure states.

  \item \textbf{Abstract interpretation:} Over-approximates program behaviour to prove absence of certain error classes, at the cost of potential false positives.
\end{itemize}

\subsection{The Tool Landscape}

Table~\ref{tab:sast_tools} summarises the major open-source SAST tools and their primary domains.

\begin{table}[htbp]
\caption{Open-Source SAST Tools Surveyed and Selected for \ACRQA{} (Tools integrated in \ACRQA{} v5.0.0rc2 shown in bold)}
\label{tab:sast_tools}
\centering
\small
\renewcommand{\arraystretch}{1.15}
\begin{tabularx}{\textwidth}{@{} l l X @{}}
\toprule
\textbf{Tool} & \textbf{Primary Language(s)} & \textbf{Focus Area} \\
\midrule
\multicolumn{3}{@{}l}{\textit{Python Language Tools}} \\
Ruff           & Python               & Fast linter; style, best-practice, and security rules \\
Bandit         & Python               & Security anti-patterns; CWE mapping; 40+ security rules \\
Semgrep        & Multi-language       & Pattern-based AST matching; custom rule sets \\
Vulture        & Python               & Unused code and dead-code detection \\
jscpd          & Multi-language       & Copy-paste and code duplication analysis \\
Radon          & Python               & Cyclomatic complexity and maintainability metrics \\
\midrule
\multicolumn{3}{@{}l}{\textit{JavaScript / TypeScript Tools}} \\
ESLint + security plugins & JavaScript / TypeScript & XSS, prototype pollution, eval misuse, injection \\
Semgrep        & Multi-language       & Shared rule sets for cross-language analysis \\
npm audit      & Node.js (deps)       & Known CVEs in Node.js dependency tree \\
\midrule
\multicolumn{3}{@{}l}{\textit{Go Language Tools}} \\
gosec          & Go                   & 35 security rules mapped to canonical IDs \\
staticcheck    & Go                   & 21 static analysis rules; style and correctness \\
\midrule
\multicolumn{3}{@{}l}{\textit{Cross-Language Scanners}} \\
Secrets Detector & All (regex-based)  & 15+ patterns: AWS keys, GitHub tokens, JWTs, DB URLs \\
SCA Scanner    & Python / Node.js     & pip-audit wrapper; CVSS$\rightarrow$severity mapping; known CVEs \\
\bottomrule
\end{tabularx}
\end{table}

Each tool excels within its niche but exhibits blind spots outside it \cite{bandit_tool,semgrep_tool,gosec_tool,eslint_security}. No single tool achieves comprehensive coverage, which motivates the multi-tool orchestration approach of \ACRQA{} \cite{ayewah2008using}.

\subsection{The False Positive Problem}

A fundamental tension in SAST design is the precision-recall trade-off. Tools configured for high recall generate large volumes of false positives --- findings that are technically triggered by a rule but do not represent real vulnerabilities in the context of the codebase. Johnson et al.\ \cite{johnson2013dont} conducted a systematic study of developer responses to static analysis warnings and found that developers frequently disable or ignore tools after encountering a sustained pattern of false positives, negating the tool's security value entirely.

Empirical studies across industrial codebases report false positive rates ranging from 35\% to over 80\% depending on tool configuration and codebase characteristics \cite{muske2016survey}. Complementary code quality metrics such as cyclomatic complexity \cite{mccabe1976complexity} and code clone density \cite{roy2009comparison} correlate with vulnerability likelihood and inform finding prioritisation. Vulnerability severity frameworks including CVSS \cite{cvss_v3} and statistical ranking approaches such as Z-ranking \cite{kremenek2003z} provide the theoretical foundation for label-free confidence scoring.

% -----------------------------------------------------------------
\section{AI-Assisted Code Review}
\label{sec:ai_code_review}

\subsection{Machine Learning and Large Language Models for Code}

Early ML approaches applied classifiers to features from code property graphs (ASTs + CFGs + PDGs) \cite{yamaguchi2014modeling} and neural models such as Code2Vec \cite{alon2019code2vec} and VulDeePecker \cite{li2018vuldeepecker}, but require large labelled datasets and generalise poorly across codebases. Transformer-based LLMs (GitHub Copilot \cite{chen2021evaluating}, CodeBERT \cite{feng2020codebert}, CodeLlama \cite{roziere2023code}) demonstrated strong zero-shot generalisation \cite{thapa2022transformer} but introduced hallucination risks: Pearce et al.\ \cite{pearce2022asleep} found Copilot generated vulnerable code in 40\% of scenarios, and Ni et al.\ \cite{ni2023chatgpt} showed detection performance is highly sensitive to prompt formulation.

\subsection{Hallucination and Consistency Challenges}

A critical limitation of LLM use in high-stakes contexts is hallucination --- the generation of plausible-sounding but factually incorrect content \cite{ji2023survey}. In a security context, a hallucinated explanation that mischaracterises exploitability or recommends an incorrect remediation can cause more harm than no explanation. Mitigations include retrieval augmentation, uncertainty quantification, and consistency-based filtering on N-gram entropy distributions \cite{kuhn2023semantic}.

% -----------------------------------------------------------------
\section{Retrieval-Augmented Generation (RAG)}
\label{sec:rag_background}

\subsection{Motivation and Architecture}

Retrieval-Augmented Generation was introduced by Lewis et al.\ \cite{lewis2020retrieval} as a method to ground LLM outputs in external factual knowledge, reducing hallucination and improving factual accuracy on knowledge-intensive tasks. The RAG architecture consists of two primary components:

\begin{enumerate}
  \item \textbf{Retriever:} Given a query, retrieves relevant documents from a pre-built knowledge base using dense or sparse vector similarity search.
  \item \textbf{Generator:} An LLM that conditions its output on both the original query and the retrieved documents, which are concatenated into the prompt context.
\end{enumerate}

The key advantage of RAG over pure parametric LLMs \cite{brown2020language} is that the knowledge base can be updated independently of the model, enabling the system to stay current with evolving security guidance without model fine-tuning.

\subsection{RAG in Security Contexts and Knowledge Base Construction}

RAG has been applied to CVE Q\&A \cite{chatsec2024}, malware analysis \cite{securityllm2024}, and automated patch generation \cite{zhang2024autocoderover}. Application to SAST finding enrichment --- as employed by \ACRQA{} --- is less explored. The quality of RAG outputs is strongly dependent on knowledge base quality \cite{guu2020retrieval}; the \ACRQA{} knowledge base of 66 security rules (422 canonical mappings) was curated from OWASP SCP \cite{owasp_scp}, NIST SP~800-53 \cite{nist_800_53}, MITRE CWE \cite{mitre_cwe}, OWASP ASVS \cite{owasp_asvs}, and OWASP API Security Top~10 \cite{owasp_api_top10}.

% -----------------------------------------------------------------
\section{Cryptographic Bill of Materials (CBoM)}
\label{sec:cbom_background}

The Software Bill of Materials (SBOM) --- a machine-readable component inventory --- gained regulatory momentum following the 2021 US Executive Order on Cybersecurity \cite{white_house_eo_2021}. Supply chain attacks have grown significantly \cite{ladisa2023sok}; SBOM formats SPDX \cite{spdx_standard} and CycloneDX \cite{cyclonedx} are now widely adopted, as is SARIF \cite{sarif_standard} for static analysis result interchange.

A Cryptographic Bill of Materials (CBoM) extends this concept to cryptographic assets --- public keys, certificates, algorithms, cipher modes --- relevant for post-quantum migration planning, as classical public-key algorithms will be vulnerable to Shor's algorithm \cite{shor1994algorithms} and Grover's algorithm \cite{grover1996fast}. NIST's PQC Standardisation project (finalised 2024) \cite{nist_pqc_2024} selected CRYSTALS-Kyber \cite{nist_fips_203}, CRYSTALS-Dilithium \cite{nist_fips_204}, and others for standardisation \cite{nist_fips_197}; organisations are recommended to begin inventory and migration immediately. The CBoM module in \ACRQA{} directly supports this by inventorying cryptographic primitives and flagging those scheduled for deprecation.

% -----------------------------------------------------------------
\section{Related Systems and Identified Gaps}
\label{sec:related_tools}

Table~\ref{tab:related_tools} summarises the principal commercial and open-source tools surveyed and their key limitations relative to \ACRQA{}.

\begin{table}[htbp]
\caption{Related Tools and Key Limitations}
\label{tab:related_tools}
\centering
\small
\renewcommand{\arraystretch}{1.25}
\begin{tabularx}{\textwidth}{@{} l X @{}}
\toprule
\textbf{Tool} & \textbf{Key Limitation vs.\ \ACRQA{}} \\
\midrule
SonarQube \cite{sonarqube}  & Enterprise licensing required for AI features; no RAG grounding \\
Snyk Code \cite{snyk}       & Cloud-dependent; proprietary AI; no CBoM or PQC awareness \\
Veracode \cite{veracode}    & Premium pricing prohibitive for academic use \\
Checkmarx \cite{checkmarx} & Enterprise cloud only; no self-hosted AI enrichment \\
DefectDojo \cite{defectdojo} & Aggregation layer only; no analysis or AI enrichment \\
Trivy \cite{trivy}          & Container/IaC focus; no cross-language SAST or LLM enrichment \\
\bottomrule
\end{tabularx}
\end{table}

From this survey, five gaps motivate \ACRQA{}: (1) no open-source platform combines 10+ SAST tools with LLM enrichment in a self-hosted service; (2) commercial AI-SAST platforms are cost-prohibitive; (3) no existing tool implements label-free confidence scoring; (4) no tool applies entropy-based LLM consistency validation; and (5) no open-source SAST tool performs CBoM generation with NIST PQC migration flagging.

% -----------------------------------------------------------------
\section{Summary}

This chapter has surveyed static analysis, AI-assisted code review, RAG, and cryptographic inventory, confirming that a self-hosted, multi-tool, AI-augmented platform with label-free confidence scoring, entropy-based consistency validation, and PQC-aware CBoM generation is a meaningful and novel contribution.

\clearpage
